\paragraph{Б.2. Отношения}
\subparagraph{Б.2.1}
Покажите, что отношение подмножества "$\subseteq$" на множестве всех 
подмножеств $\mathbb{Z}$ является отношением частичного, но не полного порядка.

\textbf{Ответ:}

Для соотвествия отношения частичному порядку, необходимо чтобы отношение 
было антисимметричным, рефлексивным и транзитивным.

1) Антисимметричным:

\[
a,b \in  2^{\left| \mathbb{Z} \right|}, a \subseteq b \text{ и } b \subseteq a \rightarrow  a = b
\]

2) рефлексивным:

\[
a \in 2^{\left| \mathbb{Z} \right|}, a \subseteq a
\]

3) транзитивным:

\[
a,b,c \in 2^{\left| \mathbb{Z} \right|}, a \subseteq b \text{ и } b \subseteq c \rightarrow  a \subseteq c
\]

Отношением частичного порядка:

Так как можно взять два множества таких что:

\[
\mathbb{X} = \{0\} \text{ и } \mathbb{Y} = \{1\} 
\]

\[
\mathbb{X} \not \subseteq \mathbb{Y} \text{ и } \mathbb{Y} \not \subseteq \mathbb{X}
\]

\subparagraph{Б.2.2}

Покажите, что для любого положительного \( n \) отношение 
"тождественности по модулю $ n $" является отношением 
эквивалентности на множестве целых чисел.  
(Мы говорим, что $ a \equiv b \pmod{n} $, 
если существует такое целое число $ q $, что $ a - b = qn $.) 
На сколько классов эквивалентности разбивает множество целых чисел это отношение?

\textbf{Ответ:}

Для соотвествия отношения эквивалентности, необходимо чтобы отношение 
было симметричным, рефлексивным и транзитивным.

1) симметричным:

Для любого $a, b \in \mathbb{N} $ можно найти $q$ такое что: $a-b=qn$,
то для $b-a=-qn$. Следовательно $aRb = bRa $

2) рефлексия:

Для $a \in \mathbb{N} $, то $a-a=0n$

3) транзитивность:

Пусть $a R b$ и $b R c$. По определению это означает, что существуют такие целые числа $q_1$ и $q_2$, что:
1. $a - b = q_1 n$
2. $b - c = q_2 n$

Нам нужно показать, что $a R c$, то есть $a - c = q_3 n$ для некоторого целого $q_3$. Сложим два предыдущих равенства:
\[ (a - b) + (b - c) = q_1 n + q_2 n \]
\[ a - c = (q_1 + q_2) n \]

Так как сумма целых чисел $(q_1 + q_2)$ — это тоже целое число (обозначим его за $q_3$), то условие выполняется. Следовательно, $a R c$.

Ответ на вторую часть вопроса:
В задаче также спрашивается: "На сколько классов эквивалентности разбивает множество целых чисел это отношение?"

\textbf{Ответ:}
Отношение разбивает множество целых чисел на $n$ классов эквивалентности. 
Эти классы соответствуют возможным остаткам от деления на $n$: $\{0, 1, 2, \dots, n-1\}$.

\subparagraph{Б.2.3}
Приведите пример отношений, которые
\begin{itemize}
    \item[\textbf{а.}] рефлексивны и симметричны, но не транзитивны;
    \item[\textbf{б.}] рефлексивны и транзитивны, но не симметричны;
    \item[\textbf{в.}] симметричны и транзитивны, но не рефлексивны.
\end{itemize}

\subparagraph{Б.2.4}
Пусть $S$ представляет собой конечное множество, а $R$ является отношением эквивалентности на $S \times S$.
Покажите, что если отношение $R$ антисимметрично, то все классы эквивалентности $S$ по отношению к $R$ содержат по одному элементу.

\subparagraph{Б.2.5}
Профессор полагает, что всякое симметричное и транзитивное отношение $R$ должно быть рефлексивным.
Он предлагает следующее доказательство: из $a R b$ в силу симметрии следует $b R a$,
а применение транзитивности к этому выводу дает $a R a$. Прав ли профессор?